% Source: https://artofproblemsolving.com/community/c3h2506266_chicken_mcnugget_theorem_for_3_or_more_numbers
% Stable identity: aops-topic-2506266-q01
\textbf{$29$ cannot be written.} If $z\ge2$ then $15z\ge30>29$, so $z\in\{0,1\}$.
\begin{itemize}
\item $z=0$: $6x+10y$ is even, but $29$ is odd.
\item $z=1$: we need $6x+10y=14$. With $y=0$ we need $6x=14$, and with $y=1$ we need $6x=4$; neither works, and $y\ge2$ gives $10y>14$.
\end{itemize}

\textbf{Every integer $n\ge30$ can be written.} The six consecutive integers $30$ to $35$ can:
\[30=6\cdot5,\quad31=6+10+15,\quad32=6\cdot2+10\cdot2,\quad33=6\cdot3+15,\quad34=6\cdot4+10,\quad35=10\cdot2+15.\]
Any $n\ge30$ is one of these plus a multiple of $6$ (subtract $6$ until the result lies in $30,\dots,35$), so it can be written too.

Hence the answer is $\boxed{29}$.

(Remark: the two-number formula does not extend in any simple way. Here every pair of $6,10,15$ has a common factor, so no pair alone can produce all large integers, yet the three together can.)
